% ------------------------------------------------------------
% Section 2 -- Related Work
% ------------------------------------------------------------
\section{Related Work}
\label{sec:related}

\paragraph{Deep LDCT denoising and its stagnation.}
RED-CNN~\cite{chen2017} remains a canonical encoder--decoder baseline for supervised LDCT denoising; adversarial approaches such as WGAN-VGG~\cite{yang2018} and DU-GAN~\cite{huang2022} trade pixel-wise fidelity for perceptual texture. The ldct-benchmark study~\cite{eulig2024} standardizes training and evaluation across this line of work and finds that, on PSNR/SSIM/VIF, ``most methods perform statistically similar and improvements over the past six years have been marginal at best.'' We adopt their exact protocol and reframe the target axis: physical fidelity rather than classical metrics. Importantly, the failure we audit is not a property of RED-CNN specifically --- it follows from the appearance-driven objectives shared by this entire family of methods; RED-CNN serves only as the standardized, representative testbed.

\paragraph{Physics-aware objectives and architectures.}
Noise-power-spectrum losses have appeared recently --- most directly CT-Mamba~\cite{li2025}, which introduces a Mamba-driven deep NPS loss, and the frequency-domain NPS loss of Zhang et al.~\cite{zhang2025}; noise-texture preservation is likewise reported as a headline property of recent dose-aware diffusion work. Closest in spirit at the architecture level, NIFA~\cite{zhao2026} embeds a noise-intensity-field formulation into the network and pairs it with a field-similarity loss to preserve realistic noise characteristics. These interventions either act purely at the loss level or embed physics into the trunk itself: none provide structural spectral guarantees (exact DC preservation of the correction, frozen near-DC gains), and none study the interaction between physics-loss weight and architectural capacity. We show empirically that loss-only physics incentives hit a weight ceiling (Sec.~\ref{sec:ceiling}) and that an architectural mechanism is required to exploit stronger incentives.

\paragraph{Frequency-adaptive restoration in natural images.}
Adaptive frequency modulation (AdaIR~\cite{cui2025} and related all-in-one restoration work) and FiLM-style conditioning~\cite{perez2018} modulate features per input, but they target natural images, operate on unconstrained feature maps, and offer no physical guarantees.

\paragraph{Gap we fill.}
Recent work thus increasingly couples LDCT denoising with noise-statistics objectives or physics-aware architectures~\cite{li2025,zhang2025,zhao2026}. To our knowledge, however, no prior LDCT work combines all four of: (i) a residual spectral head with a trunk-agnostic interface acting on the removed-noise residual; (ii) provable per-image guarantees (exact DC preservation of the residual correction, frozen near-DC gains, radial parameterization, bounded adaptive offsets); (iii) anatomy-adaptive per-image gain curves whose adaptivity is itself quantitatively measured; and (iv) an attribution protocol (matched budgets, matched loss weights, loss-ceiling control) that isolates the architectural contribution from the loss contribution. It is this specific combination --- and the isolation of the architectural contribution under matched incentives --- that we claim as novel, not any single ingredient.
